%% ECON 282E -- Session 3, Deck B2: From Loops to Arrays
%% Budget: 16 frames (~22 minutes). Provenance: slides/SOURCES.md.
%% Control flow, functions and classes are ported from QM Lec. 3 sections 5-6.
%% NumPy, broadcasting, the GrowthModel class, the vectorised Bellman update and
%% the timings are new: QM Lec. 3 section 7 is pandas/FRED, not NumPy.
%% Timings from tools/figures/s03_numbers.json, key "vectorization"; the class on
%% the later frames was run and reproduces 272 iterations and 0.00988 exactly.

%% ECON 282E -- Foundations of Macroeconomics (UC Riverside, Fall 2026)
%% Shared Beamer preamble for every deck in slides/S01 ... slides/S10.
%%
%% Usage (a deck lives in slides/S0X/ and is compiled from that folder):
%%   %% ECON 282E -- Foundations of Macroeconomics (UC Riverside, Fall 2026)
%% Shared Beamer preamble for every deck in slides/S01 ... slides/S10.
%%
%% Usage (a deck lives in slides/S0X/ and is compiled from that folder):
%%   %% ECON 282E -- Foundations of Macroeconomics (UC Riverside, Fall 2026)
%% Shared Beamer preamble for every deck in slides/S01 ... slides/S10.
%%
%% Usage (a deck lives in slides/S0X/ and is compiled from that folder):
%%   \input{../preamble/course_preamble.tex}
%%   \decktitle{1}{A1}{Who I Am \& Why This Course}{September 24, 2026}
%%   \begin{document}
%%   \makedecktitle
%%   ...
%%
%% Adapted from Courses/AI-ECON-2026/source/shared_preamble.tex (Metropolis theme,
%% concept/caution/example/takeaway boxes, \sectiondivider). Changes for this course:
%%   * course title block (\decktitle / \makedecktitle) instead of \makebeamertitle
%%   * \zh{...} typesets Chinese (book titles) under pdflatex via CJKutf8
%%   * researchbox for the "Research in the wild" frames
%%   * section pages off; TikZ libraries and the course map loaded here
%% Build with pdflatex only (two passes); tools/build_slides.py does this for you.

\documentclass[10pt,english,aspectratio=169]{beamer}

\usepackage[T1]{fontenc}
\usepackage{textcomp}
\usepackage[utf8]{inputenc}
\usepackage{babel}
\usepackage{CJKutf8}

\usepackage{amsmath, amssymb, amsthm, graphicx}
\usepackage{booktabs, multirow, tabularx}
\usepackage{tikz}
\usetikzlibrary{positioning,arrows.meta,shapes,calc,fit,backgrounds}
\usepackage{pgfplots}
\pgfplotsset{compat=1.16}
\usepackage[most]{tcolorbox}
\tcbuselibrary{skins, breakable}
\usepackage{listings}
\usepackage{xcolor}

\lstset{
  basicstyle=\ttfamily\small,
  keywordstyle=\color{blue!80!black}\bfseries,
  commentstyle=\color{green!50!black}\itshape,
  stringstyle=\color{orange!80!black},
  backgroundcolor=\color{gray!8},
  frame=single,
  rulecolor=\color{gray!40},
  breaklines=true,
  showstringspaces=false,
  numbers=none,
}

\usetheme{metropolis}
\usefonttheme{professionalfonts}
\metroset{sectionpage=none}

\definecolor{LightGray}{RGB}{230,230,230}
\definecolor{RedTitle}{RGB}{0,0,200}
\definecolor{VeryLightGray}{RGB}{245,245,245}
\definecolor{DarkText}{RGB}{0,0,0}
\definecolor{UCRBlue}{RGB}{0,45,106}
\definecolor{UCRGold}{RGB}{241,171,0}

% Stage colours of the course arc (used by the course map and elsewhere)
\colorlet{StageTrunk}{blue!12}
\colorlet{StageBranch}{cyan!14}
\colorlet{StageMachine}{gray!18}
\colorlet{StageWall}{red!14}
\colorlet{StageTools}{orange!20}
\colorlet{StageData}{green!16}
\colorlet{StageAgents}{violet!14}

\hypersetup{bookmarks=false,breaklinks=true,pdfborder={0 0 0},colorlinks=true,
  linkcolor=DarkText,urlcolor=RedTitle}

\setbeamercolor{background canvas}{bg=white}
\setbeamercolor{title}{fg=RedTitle}
\setbeamercolor{frametitle}{bg=VeryLightGray,fg=DarkText}
\setbeamercolor{normal text}{fg=DarkText}
\setbeamercolor{alerted text}{fg=RedTitle}
\setbeamersize{text margin left=5mm, text margin right=5mm}
\addtobeamertemplate{frametitle}{}{\vspace{-0.5em}}

\setbeamertemplate{navigation symbols}{}
\setbeamertemplate{footline}{%
  \hspace{5mm}{\usebeamerfont{page number in head/foot}\color{black!45}%
  ECON 282E \,$\cdot$\, \insertshortdeck}\hfill%
  \usebeamercolor[fg]{page number in head/foot}%
  \usebeamerfont{page number in head/foot}%
  \insertframenumber\,/\,\inserttotalframenumber\kern1em\vskip2pt%
}

% ---------------------------------------------------------------------------
% Chinese text under pdflatex (book titles etc.)
\newcommand{\zh}[1]{\begin{CJK*}{UTF8}{gbsn}#1\end{CJK*}}

% ---------------------------------------------------------------------------
% Deck title block
%   \decktitle{session}{deck id}{title}{date}
\newcommand{\insertshortdeck}{}
\newcommand{\decksession}{}
\newcommand{\deckid}{}
\newcommand{\decktitletext}{}
\newcommand{\deckdate}{}
\newcommand{\decktitle}[4]{%
  \renewcommand{\decksession}{#1}%
  \renewcommand{\deckid}{#2}%
  \renewcommand{\decktitletext}{#3}%
  \renewcommand{\deckdate}{#4}%
  \renewcommand{\insertshortdeck}{Session #1 \,$\cdot$\, Deck #1.#2}%
  \title{#3}%
  \author{Zhigang Feng}%
  \date{#4}%
}
\newcommand{\makedecktitle}{%
  \begin{frame}[plain,noframenumbering]
    \begin{tikzpicture}[remember picture,overlay]
      \fill[UCRBlue] (current page.north west) rectangle ([yshift=-0.55cm]current page.north east);
      \fill[UCRGold] ([yshift=-0.55cm]current page.north west) rectangle ([yshift=-0.62cm]current page.north east);
      \node[anchor=west, text=white, font=\small\bfseries] at ([xshift=5mm,yshift=-0.275cm]current page.north west)
        {ECON 282E \,$\cdot$\, Foundations of Macroeconomics};
      \node[anchor=east, text=white, font=\small] at ([xshift=-5mm,yshift=-0.275cm]current page.north east)
        {UC Riverside \,$\cdot$\, Fall 2026};
    \end{tikzpicture}
    \vspace*{1.1cm}
    {\usebeamerfont{title}\color{black!55}\large Session \decksession\ \,$\cdot$\, Deck \decksession.\deckid\par}
    \vspace{0.25cm}
    {\usebeamerfont{title}\color{RedTitle}\LARGE\bfseries \decktitletext\par}
    \vspace{0.7cm}
    {\large Zhigang Feng\par}
    \vspace{0.15cm}
    {\color{black!65}\deckdate\par}
    \vfill
    {\footnotesize\itshape\color{black!55}Build it on paper $\;\to\;$ solve it on the machine $\;\to\;$ push it to the frontier with AI\par}
  \end{frame}}

% ---------------------------------------------------------------------------
% Section divider (plain frame, not counted)
\newcommand{\sectiondivider}[2]{%
\begin{frame}[plain,noframenumbering]
\begin{center}
\vfill
{\Large\textbf{#1}\par}
\vspace{0.35cm}
{\normalsize #2\par}
\vfill
\end{center}
\end{frame}
}

% ---------------------------------------------------------------------------
% Boxes
\newtcolorbox{conceptbox}[1][]{colback=blue!3, colframe=RedTitle!55, boxrule=0.35mm,
  arc=1mm, left=2mm, right=2mm, top=1mm, bottom=1mm, #1}
\newtcolorbox{cautionbox}[1][]{colback=red!3, colframe=red!50!black, boxrule=0.35mm,
  arc=1mm, left=2mm, right=2mm, top=1mm, bottom=1mm, #1}
\newtcolorbox{examplebox}[1][]{colback=green!3, colframe=green!45!black, boxrule=0.35mm,
  arc=1mm, left=2mm, right=2mm, top=1mm, bottom=1mm, #1}
\newtcolorbox{takeawaybox}[1][]{colback=VeryLightGray, colframe=RedTitle!55, boxrule=0.35mm,
  arc=1mm, left=2mm, right=2mm, top=1mm, bottom=1mm, #1}
\newtcolorbox{researchbox}[1][]{enhanced, colback=violet!4, colframe=violet!60!black,
  boxrule=0.35mm, arc=1mm, left=2mm, right=2mm, top=1mm, bottom=1mm,
  fonttitle=\bfseries\small, title={Research in the wild}, #1}

\newcommand{\compacttables}{\renewcommand{\arraystretch}{1.15}}
\newcommand{\src}[1]{{\tiny\color{black!55}#1}}

% ---------------------------------------------------------------------------
\input{../preamble/course_map.tex}

%%   \decktitle{1}{A1}{Who I Am \& Why This Course}{September 24, 2026}
%%   \begin{document}
%%   \makedecktitle
%%   ...
%%
%% Adapted from Courses/AI-ECON-2026/source/shared_preamble.tex (Metropolis theme,
%% concept/caution/example/takeaway boxes, \sectiondivider). Changes for this course:
%%   * course title block (\decktitle / \makedecktitle) instead of \makebeamertitle
%%   * \zh{...} typesets Chinese (book titles) under pdflatex via CJKutf8
%%   * researchbox for the "Research in the wild" frames
%%   * section pages off; TikZ libraries and the course map loaded here
%% Build with pdflatex only (two passes); tools/build_slides.py does this for you.

\documentclass[10pt,english,aspectratio=169]{beamer}

\usepackage[T1]{fontenc}
\usepackage{textcomp}
\usepackage[utf8]{inputenc}
\usepackage{babel}
\usepackage{CJKutf8}

\usepackage{amsmath, amssymb, amsthm, graphicx}
\usepackage{booktabs, multirow, tabularx}
\usepackage{tikz}
\usetikzlibrary{positioning,arrows.meta,shapes,calc,fit,backgrounds}
\usepackage{pgfplots}
\pgfplotsset{compat=1.16}
\usepackage[most]{tcolorbox}
\tcbuselibrary{skins, breakable}
\usepackage{listings}
\usepackage{xcolor}

\lstset{
  basicstyle=\ttfamily\small,
  keywordstyle=\color{blue!80!black}\bfseries,
  commentstyle=\color{green!50!black}\itshape,
  stringstyle=\color{orange!80!black},
  backgroundcolor=\color{gray!8},
  frame=single,
  rulecolor=\color{gray!40},
  breaklines=true,
  showstringspaces=false,
  numbers=none,
}

\usetheme{metropolis}
\usefonttheme{professionalfonts}
\metroset{sectionpage=none}

\definecolor{LightGray}{RGB}{230,230,230}
\definecolor{RedTitle}{RGB}{0,0,200}
\definecolor{VeryLightGray}{RGB}{245,245,245}
\definecolor{DarkText}{RGB}{0,0,0}
\definecolor{UCRBlue}{RGB}{0,45,106}
\definecolor{UCRGold}{RGB}{241,171,0}

% Stage colours of the course arc (used by the course map and elsewhere)
\colorlet{StageTrunk}{blue!12}
\colorlet{StageBranch}{cyan!14}
\colorlet{StageMachine}{gray!18}
\colorlet{StageWall}{red!14}
\colorlet{StageTools}{orange!20}
\colorlet{StageData}{green!16}
\colorlet{StageAgents}{violet!14}

\hypersetup{bookmarks=false,breaklinks=true,pdfborder={0 0 0},colorlinks=true,
  linkcolor=DarkText,urlcolor=RedTitle}

\setbeamercolor{background canvas}{bg=white}
\setbeamercolor{title}{fg=RedTitle}
\setbeamercolor{frametitle}{bg=VeryLightGray,fg=DarkText}
\setbeamercolor{normal text}{fg=DarkText}
\setbeamercolor{alerted text}{fg=RedTitle}
\setbeamersize{text margin left=5mm, text margin right=5mm}
\addtobeamertemplate{frametitle}{}{\vspace{-0.5em}}

\setbeamertemplate{navigation symbols}{}
\setbeamertemplate{footline}{%
  \hspace{5mm}{\usebeamerfont{page number in head/foot}\color{black!45}%
  ECON 282E \,$\cdot$\, \insertshortdeck}\hfill%
  \usebeamercolor[fg]{page number in head/foot}%
  \usebeamerfont{page number in head/foot}%
  \insertframenumber\,/\,\inserttotalframenumber\kern1em\vskip2pt%
}

% ---------------------------------------------------------------------------
% Chinese text under pdflatex (book titles etc.)
\newcommand{\zh}[1]{\begin{CJK*}{UTF8}{gbsn}#1\end{CJK*}}

% ---------------------------------------------------------------------------
% Deck title block
%   \decktitle{session}{deck id}{title}{date}
\newcommand{\insertshortdeck}{}
\newcommand{\decksession}{}
\newcommand{\deckid}{}
\newcommand{\decktitletext}{}
\newcommand{\deckdate}{}
\newcommand{\decktitle}[4]{%
  \renewcommand{\decksession}{#1}%
  \renewcommand{\deckid}{#2}%
  \renewcommand{\decktitletext}{#3}%
  \renewcommand{\deckdate}{#4}%
  \renewcommand{\insertshortdeck}{Session #1 \,$\cdot$\, Deck #1.#2}%
  \title{#3}%
  \author{Zhigang Feng}%
  \date{#4}%
}
\newcommand{\makedecktitle}{%
  \begin{frame}[plain,noframenumbering]
    \begin{tikzpicture}[remember picture,overlay]
      \fill[UCRBlue] (current page.north west) rectangle ([yshift=-0.55cm]current page.north east);
      \fill[UCRGold] ([yshift=-0.55cm]current page.north west) rectangle ([yshift=-0.62cm]current page.north east);
      \node[anchor=west, text=white, font=\small\bfseries] at ([xshift=5mm,yshift=-0.275cm]current page.north west)
        {ECON 282E \,$\cdot$\, Foundations of Macroeconomics};
      \node[anchor=east, text=white, font=\small] at ([xshift=-5mm,yshift=-0.275cm]current page.north east)
        {UC Riverside \,$\cdot$\, Fall 2026};
    \end{tikzpicture}
    \vspace*{1.1cm}
    {\usebeamerfont{title}\color{black!55}\large Session \decksession\ \,$\cdot$\, Deck \decksession.\deckid\par}
    \vspace{0.25cm}
    {\usebeamerfont{title}\color{RedTitle}\LARGE\bfseries \decktitletext\par}
    \vspace{0.7cm}
    {\large Zhigang Feng\par}
    \vspace{0.15cm}
    {\color{black!65}\deckdate\par}
    \vfill
    {\footnotesize\itshape\color{black!55}Build it on paper $\;\to\;$ solve it on the machine $\;\to\;$ push it to the frontier with AI\par}
  \end{frame}}

% ---------------------------------------------------------------------------
% Section divider (plain frame, not counted)
\newcommand{\sectiondivider}[2]{%
\begin{frame}[plain,noframenumbering]
\begin{center}
\vfill
{\Large\textbf{#1}\par}
\vspace{0.35cm}
{\normalsize #2\par}
\vfill
\end{center}
\end{frame}
}

% ---------------------------------------------------------------------------
% Boxes
\newtcolorbox{conceptbox}[1][]{colback=blue!3, colframe=RedTitle!55, boxrule=0.35mm,
  arc=1mm, left=2mm, right=2mm, top=1mm, bottom=1mm, #1}
\newtcolorbox{cautionbox}[1][]{colback=red!3, colframe=red!50!black, boxrule=0.35mm,
  arc=1mm, left=2mm, right=2mm, top=1mm, bottom=1mm, #1}
\newtcolorbox{examplebox}[1][]{colback=green!3, colframe=green!45!black, boxrule=0.35mm,
  arc=1mm, left=2mm, right=2mm, top=1mm, bottom=1mm, #1}
\newtcolorbox{takeawaybox}[1][]{colback=VeryLightGray, colframe=RedTitle!55, boxrule=0.35mm,
  arc=1mm, left=2mm, right=2mm, top=1mm, bottom=1mm, #1}
\newtcolorbox{researchbox}[1][]{enhanced, colback=violet!4, colframe=violet!60!black,
  boxrule=0.35mm, arc=1mm, left=2mm, right=2mm, top=1mm, bottom=1mm,
  fonttitle=\bfseries\small, title={Research in the wild}, #1}

\newcommand{\compacttables}{\renewcommand{\arraystretch}{1.15}}
\newcommand{\src}[1]{{\tiny\color{black!55}#1}}

% ---------------------------------------------------------------------------
%% Course map: the recurring "you are here" slide for ECON 282E.
%% Loaded by course_preamble.tex. Pattern follows AI-ECON-2026 lec01_map.tex, but the map
%% shows the whole ten-session course rather than one lecture.
%%
%%   \CourseMap{s}{label}   s = 1..10 highlights session s and prints "you are here: label"
%%                          (e.g. \CourseMap{1}{1.A2}); earlier sessions get a check mark.
%%   \CourseMap{0}{}        full overview, nothing highlighted.
%%
%% Note: TikZ option lists are comma-split before TeX conditionals run, so every \ifnum
%% below sits outside [...] option lists.

\newcommand{\CourseMap}[2]{%
\begin{center}
\begin{tikzpicture}[
    sess/.style={rectangle, rounded corners=2pt, draw=black!45, minimum width=1.34cm,
        minimum height=1.12cm, text width=1.26cm, align=center, inner sep=1pt},
    stage/.style={font=\tiny\bfseries, text=black!70},
]
  \foreach \i/\d/\t/\c in {%
      1/Sep 24/Intro \\ Growth/StageTrunk,
      2/Oct 1/HA $\cdot$ OLG \\ Policy/StageBranch,
      3/Oct 8/Num.\ DP \\ Python/StageMachine,
      4/Oct 15/Numerics \\ I/StageMachine,
      5/Oct 22/Numerics \\ II/StageMachine,
      6/Oct 29/The wall \\ \textbf{Midterm}/StageWall,
      7/Nov 5/ML \\ Deep nets/StageTools,
      8/Nov 12/RL \\ HA $+$ DL/StageTools,
      9/Nov 19/Text \\ LLMs/StageData,
      10/Dec 3/Agents \\ \textbf{Projects}/StageAgents} {
    \node[sess, fill=\c] (n\i) at ({1.46*(\i-1)}, 0)
      {{\scriptsize\bfseries S\i}\\[-1pt]{\tiny\color{black!60}\d}\\[1pt]{\tiny \t}};
    \ifnum\i<#1
      \node[font=\scriptsize, text=green!45!black] at ([xshift=-2.5mm,yshift=-2.2mm]n\i.north east) {$\checkmark$};
    \fi
    \ifnum\i=#1
      \draw[RedTitle, very thick, rounded corners=2pt]
        ([xshift=-0.6mm,yshift=-0.6mm]n\i.south west) rectangle ([xshift=0.6mm,yshift=0.6mm]n\i.north east);
      % keep the label inside the picture at both ends of the row
      \ifnum\i<3
        \node[font=\scriptsize\bfseries, text=RedTitle, anchor=south west] at ([yshift=1.2mm]n\i.north west)
          {$\blacktriangledown$\,you are here: #2};
      \else\ifnum\i>8
        \node[font=\scriptsize\bfseries, text=RedTitle, anchor=south east] at ([yshift=1.2mm]n\i.north east)
          {you are here: #2\,$\blacktriangledown$};
      \else
        \node[font=\scriptsize\bfseries, text=RedTitle, anchor=south] at ([yshift=1.2mm]n\i.north)
          {you are here: #2\,$\blacktriangledown$};
      \fi\fi
    \fi
  }
  % stage band under the sessions
  \foreach \a/\b/\lab/\c in {1/1/trunk/StageTrunk, 2/2/branches/StageBranch,
      3/5/the machine $\cdot$ classical methods/StageMachine, 6/6/the wall/StageWall,
      7/8/new tools/StageTools, 9/9/new data/StageData, 10/10/colleagues/StageAgents} {
    \fill[\c] ([yshift=-1.5mm]n\a.south west) rectangle ([yshift=-4.6mm]n\b.south east);
    \path ([yshift=-1.5mm]n\a.south west) -- ([yshift=-4.6mm]n\b.south east)
      node[midway, stage] {\lab};
  }
  \node[font=\footnotesize\itshape, text=black!70] at ([yshift=-9.5mm]$(n1.south)!0.5!(n10.south)$)
    {Build it on paper $\;\to\;$ solve it on the machine $\;\to\;$ push it to the frontier with AI};
\end{tikzpicture}
\end{center}
}


%%   \decktitle{1}{A1}{Who I Am \& Why This Course}{September 24, 2026}
%%   \begin{document}
%%   \makedecktitle
%%   ...
%%
%% Adapted from Courses/AI-ECON-2026/source/shared_preamble.tex (Metropolis theme,
%% concept/caution/example/takeaway boxes, \sectiondivider). Changes for this course:
%%   * course title block (\decktitle / \makedecktitle) instead of \makebeamertitle
%%   * \zh{...} typesets Chinese (book titles) under pdflatex via CJKutf8
%%   * researchbox for the "Research in the wild" frames
%%   * section pages off; TikZ libraries and the course map loaded here
%% Build with pdflatex only (two passes); tools/build_slides.py does this for you.

\documentclass[10pt,english,aspectratio=169]{beamer}

\usepackage[T1]{fontenc}
\usepackage{textcomp}
\usepackage[utf8]{inputenc}
\usepackage{babel}
\usepackage{CJKutf8}

\usepackage{amsmath, amssymb, amsthm, graphicx}
\usepackage{booktabs, multirow, tabularx}
\usepackage{tikz}
\usetikzlibrary{positioning,arrows.meta,shapes,calc,fit,backgrounds}
\usepackage{pgfplots}
\pgfplotsset{compat=1.16}
\usepackage[most]{tcolorbox}
\tcbuselibrary{skins, breakable}
\usepackage{listings}
\usepackage{xcolor}

\lstset{
  basicstyle=\ttfamily\small,
  keywordstyle=\color{blue!80!black}\bfseries,
  commentstyle=\color{green!50!black}\itshape,
  stringstyle=\color{orange!80!black},
  backgroundcolor=\color{gray!8},
  frame=single,
  rulecolor=\color{gray!40},
  breaklines=true,
  showstringspaces=false,
  numbers=none,
}

\usetheme{metropolis}
\usefonttheme{professionalfonts}
\metroset{sectionpage=none}

\definecolor{LightGray}{RGB}{230,230,230}
\definecolor{RedTitle}{RGB}{0,0,200}
\definecolor{VeryLightGray}{RGB}{245,245,245}
\definecolor{DarkText}{RGB}{0,0,0}
\definecolor{UCRBlue}{RGB}{0,45,106}
\definecolor{UCRGold}{RGB}{241,171,0}

% Stage colours of the course arc (used by the course map and elsewhere)
\colorlet{StageTrunk}{blue!12}
\colorlet{StageBranch}{cyan!14}
\colorlet{StageMachine}{gray!18}
\colorlet{StageWall}{red!14}
\colorlet{StageTools}{orange!20}
\colorlet{StageData}{green!16}
\colorlet{StageAgents}{violet!14}

\hypersetup{bookmarks=false,breaklinks=true,pdfborder={0 0 0},colorlinks=true,
  linkcolor=DarkText,urlcolor=RedTitle}

\setbeamercolor{background canvas}{bg=white}
\setbeamercolor{title}{fg=RedTitle}
\setbeamercolor{frametitle}{bg=VeryLightGray,fg=DarkText}
\setbeamercolor{normal text}{fg=DarkText}
\setbeamercolor{alerted text}{fg=RedTitle}
\setbeamersize{text margin left=5mm, text margin right=5mm}
\addtobeamertemplate{frametitle}{}{\vspace{-0.5em}}

\setbeamertemplate{navigation symbols}{}
\setbeamertemplate{footline}{%
  \hspace{5mm}{\usebeamerfont{page number in head/foot}\color{black!45}%
  ECON 282E \,$\cdot$\, \insertshortdeck}\hfill%
  \usebeamercolor[fg]{page number in head/foot}%
  \usebeamerfont{page number in head/foot}%
  \insertframenumber\,/\,\inserttotalframenumber\kern1em\vskip2pt%
}

% ---------------------------------------------------------------------------
% Chinese text under pdflatex (book titles etc.)
\newcommand{\zh}[1]{\begin{CJK*}{UTF8}{gbsn}#1\end{CJK*}}

% ---------------------------------------------------------------------------
% Deck title block
%   \decktitle{session}{deck id}{title}{date}
\newcommand{\insertshortdeck}{}
\newcommand{\decksession}{}
\newcommand{\deckid}{}
\newcommand{\decktitletext}{}
\newcommand{\deckdate}{}
\newcommand{\decktitle}[4]{%
  \renewcommand{\decksession}{#1}%
  \renewcommand{\deckid}{#2}%
  \renewcommand{\decktitletext}{#3}%
  \renewcommand{\deckdate}{#4}%
  \renewcommand{\insertshortdeck}{Session #1 \,$\cdot$\, Deck #1.#2}%
  \title{#3}%
  \author{Zhigang Feng}%
  \date{#4}%
}
\newcommand{\makedecktitle}{%
  \begin{frame}[plain,noframenumbering]
    \begin{tikzpicture}[remember picture,overlay]
      \fill[UCRBlue] (current page.north west) rectangle ([yshift=-0.55cm]current page.north east);
      \fill[UCRGold] ([yshift=-0.55cm]current page.north west) rectangle ([yshift=-0.62cm]current page.north east);
      \node[anchor=west, text=white, font=\small\bfseries] at ([xshift=5mm,yshift=-0.275cm]current page.north west)
        {ECON 282E \,$\cdot$\, Foundations of Macroeconomics};
      \node[anchor=east, text=white, font=\small] at ([xshift=-5mm,yshift=-0.275cm]current page.north east)
        {UC Riverside \,$\cdot$\, Fall 2026};
    \end{tikzpicture}
    \vspace*{1.1cm}
    {\usebeamerfont{title}\color{black!55}\large Session \decksession\ \,$\cdot$\, Deck \decksession.\deckid\par}
    \vspace{0.25cm}
    {\usebeamerfont{title}\color{RedTitle}\LARGE\bfseries \decktitletext\par}
    \vspace{0.7cm}
    {\large Zhigang Feng\par}
    \vspace{0.15cm}
    {\color{black!65}\deckdate\par}
    \vfill
    {\footnotesize\itshape\color{black!55}Build it on paper $\;\to\;$ solve it on the machine $\;\to\;$ push it to the frontier with AI\par}
  \end{frame}}

% ---------------------------------------------------------------------------
% Section divider (plain frame, not counted)
\newcommand{\sectiondivider}[2]{%
\begin{frame}[plain,noframenumbering]
\begin{center}
\vfill
{\Large\textbf{#1}\par}
\vspace{0.35cm}
{\normalsize #2\par}
\vfill
\end{center}
\end{frame}
}

% ---------------------------------------------------------------------------
% Boxes
\newtcolorbox{conceptbox}[1][]{colback=blue!3, colframe=RedTitle!55, boxrule=0.35mm,
  arc=1mm, left=2mm, right=2mm, top=1mm, bottom=1mm, #1}
\newtcolorbox{cautionbox}[1][]{colback=red!3, colframe=red!50!black, boxrule=0.35mm,
  arc=1mm, left=2mm, right=2mm, top=1mm, bottom=1mm, #1}
\newtcolorbox{examplebox}[1][]{colback=green!3, colframe=green!45!black, boxrule=0.35mm,
  arc=1mm, left=2mm, right=2mm, top=1mm, bottom=1mm, #1}
\newtcolorbox{takeawaybox}[1][]{colback=VeryLightGray, colframe=RedTitle!55, boxrule=0.35mm,
  arc=1mm, left=2mm, right=2mm, top=1mm, bottom=1mm, #1}
\newtcolorbox{researchbox}[1][]{enhanced, colback=violet!4, colframe=violet!60!black,
  boxrule=0.35mm, arc=1mm, left=2mm, right=2mm, top=1mm, bottom=1mm,
  fonttitle=\bfseries\small, title={Research in the wild}, #1}

\newcommand{\compacttables}{\renewcommand{\arraystretch}{1.15}}
\newcommand{\src}[1]{{\tiny\color{black!55}#1}}

% ---------------------------------------------------------------------------
%% Course map: the recurring "you are here" slide for ECON 282E.
%% Loaded by course_preamble.tex. Pattern follows AI-ECON-2026 lec01_map.tex, but the map
%% shows the whole ten-session course rather than one lecture.
%%
%%   \CourseMap{s}{label}   s = 1..10 highlights session s and prints "you are here: label"
%%                          (e.g. \CourseMap{1}{1.A2}); earlier sessions get a check mark.
%%   \CourseMap{0}{}        full overview, nothing highlighted.
%%
%% Note: TikZ option lists are comma-split before TeX conditionals run, so every \ifnum
%% below sits outside [...] option lists.

\newcommand{\CourseMap}[2]{%
\begin{center}
\begin{tikzpicture}[
    sess/.style={rectangle, rounded corners=2pt, draw=black!45, minimum width=1.34cm,
        minimum height=1.12cm, text width=1.26cm, align=center, inner sep=1pt},
    stage/.style={font=\tiny\bfseries, text=black!70},
]
  \foreach \i/\d/\t/\c in {%
      1/Sep 24/Intro \\ Growth/StageTrunk,
      2/Oct 1/HA $\cdot$ OLG \\ Policy/StageBranch,
      3/Oct 8/Num.\ DP \\ Python/StageMachine,
      4/Oct 15/Numerics \\ I/StageMachine,
      5/Oct 22/Numerics \\ II/StageMachine,
      6/Oct 29/The wall \\ \textbf{Midterm}/StageWall,
      7/Nov 5/ML \\ Deep nets/StageTools,
      8/Nov 12/RL \\ HA $+$ DL/StageTools,
      9/Nov 19/Text \\ LLMs/StageData,
      10/Dec 3/Agents \\ \textbf{Projects}/StageAgents} {
    \node[sess, fill=\c] (n\i) at ({1.46*(\i-1)}, 0)
      {{\scriptsize\bfseries S\i}\\[-1pt]{\tiny\color{black!60}\d}\\[1pt]{\tiny \t}};
    \ifnum\i<#1
      \node[font=\scriptsize, text=green!45!black] at ([xshift=-2.5mm,yshift=-2.2mm]n\i.north east) {$\checkmark$};
    \fi
    \ifnum\i=#1
      \draw[RedTitle, very thick, rounded corners=2pt]
        ([xshift=-0.6mm,yshift=-0.6mm]n\i.south west) rectangle ([xshift=0.6mm,yshift=0.6mm]n\i.north east);
      % keep the label inside the picture at both ends of the row
      \ifnum\i<3
        \node[font=\scriptsize\bfseries, text=RedTitle, anchor=south west] at ([yshift=1.2mm]n\i.north west)
          {$\blacktriangledown$\,you are here: #2};
      \else\ifnum\i>8
        \node[font=\scriptsize\bfseries, text=RedTitle, anchor=south east] at ([yshift=1.2mm]n\i.north east)
          {you are here: #2\,$\blacktriangledown$};
      \else
        \node[font=\scriptsize\bfseries, text=RedTitle, anchor=south] at ([yshift=1.2mm]n\i.north)
          {you are here: #2\,$\blacktriangledown$};
      \fi\fi
    \fi
  }
  % stage band under the sessions
  \foreach \a/\b/\lab/\c in {1/1/trunk/StageTrunk, 2/2/branches/StageBranch,
      3/5/the machine $\cdot$ classical methods/StageMachine, 6/6/the wall/StageWall,
      7/8/new tools/StageTools, 9/9/new data/StageData, 10/10/colleagues/StageAgents} {
    \fill[\c] ([yshift=-1.5mm]n\a.south west) rectangle ([yshift=-4.6mm]n\b.south east);
    \path ([yshift=-1.5mm]n\a.south west) -- ([yshift=-4.6mm]n\b.south east)
      node[midway, stage] {\lab};
  }
  \node[font=\footnotesize\itshape, text=black!70] at ([yshift=-9.5mm]$(n1.south)!0.5!(n10.south)$)
    {Build it on paper $\;\to\;$ solve it on the machine $\;\to\;$ push it to the frontier with AI};
\end{tikzpicture}
\end{center}
}


\graphicspath{{./figures/}}

\lstset{language=Python, basicstyle=\ttfamily\scriptsize, frame=none,
        backgroundcolor=\color{gray!8}, aboveskip=2pt, belowskip=2pt}

\decktitle{3}{B2}{From Loops to Arrays}{Thursday, October 8, 2026}

\begin{document}

\makedecktitle

%=============================================================================
\begin{frame}{Where We Are}
\CourseMap{3}{3.B2}
\vspace{-1mm}
\begin{center}
\small \textbf{Block B:} \textbf{3.B1} your toolkit $\;\cdot\;$ \textbf{3.B2} from loops to arrays $\;\cdot\;$ \textbf{3.B3} PyTorch, and coding with AI $\;\cdot\;$ \textbf{3.B4} the data workflow.
\end{center}
\end{frame}

%=============================================================================
\begin{frame}[fragile]{One Bellman Update, Written Three Ways}
\begin{columns}[T]
\begin{column}{0.53\textwidth}
\footnotesize
\textbf{1. On the board}
\[
  v_{n+1}(k_i)=\max_{j}\big\{U_{ij}+\beta v_n(k_j)\big\}
\]
\vspace{-3mm}
\begin{lstlisting}
# 2. As you would say it out loud
for i in range(N):
    best = -1e18
    for j in range(N):
        val = U[i, j] + beta*V[j]
        if val > best:
            best = val
    Vn[i] = best

# 3. As NumPy wants to hear it
Vn = (U + beta*V[None, :]).max(axis=1)
\end{lstlisting}
\end{column}
\begin{column}{0.44\textwidth}
\small
Identical arithmetic. Identical answer --- \texttt{max|difference| = 0}, not ``close''.
\medskip
\begin{center}\footnotesize
\begin{tabular}{@{}lrrr@{}}
\toprule
$N$ & loop & array & speed-up \\
\midrule
200 & 0.042\,s & 0.0007\,s & $64\times$ \\
500 & 0.263\,s & 0.0039\,s & $68\times$ \\
1000 & 1.047\,s & 0.0159\,s & $\mathbf{66\times}$ \\
\bottomrule
\end{tabular}
\end{center}
\vspace{1mm}
\begin{conceptbox}\footnotesize
Multiply the last row by the 272 sweeps of 3.A2: \textbf{4.7 minutes against 4.3 seconds}, for the same number.\\[1mm]
{\scriptsize Timings move a few percent run to run; the ratio does not.}
\end{conceptbox}
\end{column}
\end{columns}
\end{frame}

%=============================================================================
\begin{frame}[fragile]{Control Flow, and Where It Belongs}
\begin{columns}[T]
\begin{column}{0.51\textwidth}
\begin{lstlisting}
if r > rho:
    print("over-accumulation")
elif r < 0:
    print("negative real rate")
else:
    print("fine")

for n in range(max_iter):     # known count
    ...
while err > tol:              # unknown count
    ...

# a comprehension: shorter, still a loop
booms = [x for x in data if x > 0]
\end{lstlisting}
\end{column}
\begin{column}{0.46\textwidth}
\small
Indentation \emph{is} the block structure --- there are no braces, and a stray space is a syntax error.
\medskip
\begin{takeawaybox}[title=\textbf{The rule for this course}]\footnotesize
\texttt{while} belongs on the \textbf{outer} loop, where the iteration count is genuinely unknown: ``keep sweeping until converged''.\\[1mm]
The \textbf{inner} loops over states and choices should not be Python loops at all.
\end{takeawaybox}
\medskip
\begin{cautionbox}\footnotesize
A comprehension reads as set-builder notation, $\{x : x\in\text{data},\,x>0\}$ --- but it is \textbf{shorter, not faster}. Use it for structure, arrays for arithmetic.
\end{cautionbox}
\end{column}
\end{columns}
\vspace{1mm}
\src{Adapted from QM Lec.~3, ``Conditionals and Loops'' and ``List Comprehensions''.}
\end{frame}

%=============================================================================
\begin{frame}[fragile]{Functions: Defaults, Docstrings, and Returning More Than One Thing}
\begin{columns}[T]
\begin{column}{0.54\textwidth}
\begin{lstlisting}
def cobb_douglas(k, z=1.0, alpha=0.36):
    """Output per worker, y = z k^alpha."""
    return z * k**alpha

cobb_douglas(2.0)                # 1.319
cobb_douglas(2.0, alpha=0.40)    # named

def solve(model, tol=1e-6):
    ...
    return V, g, n               # a tuple

V, g, n = solve(m)               # unpacked
\end{lstlisting}
\end{column}
\begin{column}{0.43\textwidth}
\small
\begin{itemize}\footnotesize
  \item Default arguments put the \textbf{calibration} in one visible place.
  \item Named arguments at the call site make \texttt{solve(m, 1e-9)} readable again.
  \item Returning a tuple is how one function hands back a value, a policy and a diagnostic.
\end{itemize}
\medskip
\begin{cautionbox}\footnotesize
Never use a \textbf{mutable} default (\texttt{def f(x, hist=[])}). It is created once and shared by every call --- the reference-versus-copy trap of 3.B1 in its nastiest form.
\end{cautionbox}
\end{column}
\end{columns}
\vspace{1mm}
\src{Adapted from QM Lec.~3, ``Functions: Reusable Blocks of Logic''.}
\end{frame}

%=============================================================================
\begin{frame}[fragile]{Modules: When One File Stops Being Enough}
\begin{columns}[T]
\begin{column}{0.50\textwidth}
\begin{lstlisting}
import numpy as np
from scipy.optimize import brentq

# your own file, growth.py
from growth import GrowthModel

if __name__ == "__main__":
    # runs only when executed directly,
    # not when imported
    main()
\end{lstlisting}
\end{column}
\begin{column}{0.47\textwidth}
\small
A module is just a \texttt{.py} file. Splitting one gives you the project shape that survives to Session 5:
\medskip
\begin{center}\footnotesize
\begin{tabular}{@{}ll@{}}
\toprule
\texttt{growth.py} & the model and its methods \\
\texttt{solvers.py} & VFI, time iteration \\
\texttt{run.py} & calibration, output \\
\bottomrule
\end{tabular}
\end{center}
\medskip
\begin{examplebox}\footnotesize
The \texttt{\_\_main\_\_} guard is what lets \texttt{run.py} be both a script you submit to a scheduler and a module a notebook imports.
\end{examplebox}
\end{column}
\end{columns}
\vspace{1mm}
\src{Adapted from QM Lec.~3, ``Modules and Imports''.}
\end{frame}

%=============================================================================
\begin{frame}{Why a Model Wants to Be a Class}
\begin{columns}[T]
\begin{column}{0.49\textwidth}
\small
A model is \textbf{state} plus \textbf{behaviour}. The state is $\alpha,\beta,\delta$, the grid, the chain; the behaviour is ``take one Bellman step'', ``solve'', ``report the exact answer''.
\medskip

Passing eleven loose variables into every function is how the wrong $\beta$ ends up in the accuracy check. A class ties them together so that cannot happen.
\end{column}
\begin{column}{0.48\textwidth}
\begin{center}\footnotesize
\begin{tabularx}{\linewidth}{@{}>{\raggedright\arraybackslash}p{2.1cm} >{\raggedright\arraybackslash}X@{}}
\toprule
\textbf{Without} & \texttt{solve(k, U, beta, tol, alpha, z, Pi, ...)} \\
\textbf{With} & \texttt{m.solve(tol)} \\
\bottomrule
\end{tabularx}
\end{center}
\medskip
\begin{takeawaybox}\footnotesize
And it is the same pattern you will meet again in 3.B3: a PyTorch network \emph{is} a class, holding parameters as state and a \texttt{forward} method as behaviour.
\end{takeawaybox}
\end{column}
\end{columns}
\vspace{1mm}
\src{Adapted from QM Lec.~3, ``Why Economists Need OOP'' and ``Without Classes vs With Classes''.}
\end{frame}

%=============================================================================
\begin{frame}[fragile]{Anatomy of a Class: Setting Up the Model}
\begin{columns}[T]
\begin{column}{0.56\textwidth}
\begin{lstlisting}
class GrowthModel:
    """u(c)=ln c, y = z k^alpha, delta = 1."""

    def __init__(self, alpha=0.36, beta=0.95,
                 n=200, k_lo=0.05, k_hi=0.60):
        self.alpha, self.beta = alpha, beta
        self.k = np.linspace(k_lo, k_hi, n)

        C = (self.k**alpha)[:, None] \
            - self.k[None, :]
        self.U = np.full_like(C, -1e10)
        np.log(C, out=self.U, where=C > 0)
\end{lstlisting}
\end{column}
\begin{column}{0.41\textwidth}
\small
\texttt{\_\_init\_\_} runs once, when the object is made. Everything it assigns to \texttt{self} is available to every method afterwards.
\medskip
\begin{conceptbox}\footnotesize
Note \emph{what} it precomputes: the payoff matrix $U$, which does not depend on $v$. Building it in the constructor is the same optimisation as 3.A2's ``build it once, reuse it 272 times'' --- now enforced by the structure.
\end{conceptbox}
\end{column}
\end{columns}
\vspace{1mm}
\src{Class syntax adapted from QM Lec.~3, ``Anatomy of a Class''; \texttt{GrowthModel} written for this course.}
\end{frame}

%=============================================================================
\begin{frame}[fragile]{The Same Class, Now with Behaviour}
\begin{columns}[T]
\begin{column}{0.54\textwidth}
\begin{lstlisting}
    def bellman(self, V):
        M = self.U + self.beta * V[None, :]
        return M.max(axis=1), M.argmax(axis=1)

    def solve(self, tol=1e-6):
        V, n = np.zeros(self.k.size), 0
        while True:
            Vn, pol = self.bellman(V)
            d = np.max(np.abs(Vn - V))
            V, n = Vn, n + 1
            if d < tol:
                return V, self.k[pol], n

    def exact_policy(self):
        return self.alpha*self.beta*self.k**self.alpha
\end{lstlisting}
\end{column}
\begin{column}{0.43\textwidth}
\begin{lstlisting}
m = GrowthModel()
V, g, n = m.solve()
n
# 272
np.max(np.abs(g/m.exact_policy() - 1))
# 0.009879
\end{lstlisting}
\vspace{-1mm}
\begin{takeawaybox}\footnotesize
Same two numbers as this morning, to every digit. Restructuring code must not move a result --- if it does, one of the two versions was wrong.
\end{takeawaybox}
\medskip
\begin{examplebox}\footnotesize
\texttt{exact\_policy} lives \emph{inside} the model, so the benchmark can never be evaluated at the wrong $\alpha$.
\end{examplebox}
\end{column}
\end{columns}
\end{frame}

%=============================================================================
\begin{frame}[fragile]{NumPy: One Type, One Block of Memory}
\begin{columns}[T]
\begin{column}{0.50\textwidth}
\begin{lstlisting}
k = np.linspace(0.05, 0.6, 200)
k.shape, k.dtype     # (200,) float64
k.nbytes             # 1600

A = np.zeros((200, 200))
A.shape              # (200, 200)
A.nbytes             # 320000

s = k[10:20]         # a VIEW
s[0] = 99.0          # changes k!
c = k[10:20].copy()  # a copy
\end{lstlisting}
\end{column}
\begin{column}{0.47\textwidth}
\small
A Python list stores pointers to objects scattered in memory. A NumPy array stores \textbf{raw numbers, contiguously, all of one type} --- which is what lets compiled code march through them.
\medskip
\begin{cautionbox}\footnotesize
A slice is a \textbf{view}, not a copy: 3.B1's reference-versus-copy rule, now with teeth. Writing into \texttt{V[10:20]} writes into \texttt{V}.
\end{cautionbox}
\medskip
\begin{examplebox}\footnotesize
\texttt{dtype} is where 3.A1's float64 decision actually gets made --- and \texttt{np.float32} silently halves your precision.
\end{examplebox}
\end{column}
\end{columns}
\end{frame}

%=============================================================================
\begin{frame}[fragile]{Broadcasting: The Idea Behind the Whole Speed-Up}
\begin{columns}[T]
\begin{column}{0.49\textwidth}
\small
Two arrays are compatible if, reading shapes from the right, each pair of dimensions is \textbf{equal} or one of them is \textbf{1}. The size-1 dimension is then stretched --- without ever being stored.
\medskip
\begin{lstlisting}
k[:, None].shape     # (200, 1)  rows
k[None, :].shape     # (1, 200)  columns

C = (k**alpha)[:, None] - k[None, :]
C.shape              # (200, 200)
# C[i, j] = k_i^alpha - k_j
\end{lstlisting}
\end{column}
\begin{column}{0.48\textwidth}
\centering
\begin{tikzpicture}[font=\tiny]
  \node[anchor=north west, align=left] at (0,0)
    {\texttt{(200, 1)}\\ \texttt{(\ \ 1,200)}\\[1pt] \texttt{---------}\\ \texttt{(200,200)}};
  \draw[RedTitle, thick, ->] (2.3,-0.18) -- (3.1,-0.18);
  \node[anchor=west, font=\tiny, text=RedTitle] at (3.15,-0.18) {stretched};
  \draw[RedTitle, thick, ->] (2.3,-0.46) -- (3.1,-0.46);
  \node[anchor=west, font=\tiny, text=RedTitle] at (3.15,-0.46) {stretched};
\end{tikzpicture}
\\[2mm]
\begin{conceptbox}\footnotesize
\textbf{Every state against every choice, in one expression.} That is the double loop of 3.A2, expressed as a shape.
\end{conceptbox}
\medskip
\begin{cautionbox}\footnotesize
The danger is that broadcasting almost never errors. Get an axis wrong and you get a plausible array of the wrong shape --- so \textbf{assert the shape}.
\end{cautionbox}
\end{column}
\end{columns}
\end{frame}

%=============================================================================
\begin{frame}[fragile]{The Vectorised Bellman Update, Line by Line}
\begin{columns}[T]
\begin{column}{0.52\textwidth}
\begin{lstlisting}
M  = U + beta * V[None, :]
#    (N,N) + (1,N) -> (N,N)
#    M[i, j] = U[i,j] + beta*V[j]

Vn = M.max(axis=1)      # over choices j
pol = M.argmax(axis=1)  # which j won
\end{lstlisting}
\vspace{-1mm}
\small
With a shock the same idea carries one more axis:
\vspace{-1mm}
\begin{lstlisting}
EV = V @ Pi.T                # (N, nz)
M  = U + beta * EV.T[None]   # (N,nz,N)
Vn = M.max(axis=2)
\end{lstlisting}
\end{column}
\begin{column}{0.45\textwidth}
\begin{takeawaybox}[title=\textbf{The recipe}]\small
Give every index of the mathematics \textbf{its own axis}, let broadcasting form the full object, and reduce along the axis you were maximising over.
\end{takeawaybox}
\medskip
\begin{examplebox}\footnotesize
The expectation is a \textbf{matrix product}, \texttt{V @ Pi.T}, done once per sweep --- not re-summed inside the maximization, which is where the naive version wastes most of its time.
\end{examplebox}
\end{column}
\end{columns}
\end{frame}

%=============================================================================
\begin{frame}{What Vectorising Does Not Buy You}
\begin{columns}[T]
\begin{column}{0.51\textwidth}
\small
The trick is to build the whole $(N,N)$ object at once. So the memory bill is $N^2$ doubles, and it arrives all at once:
\medskip
\begin{center}\footnotesize
\begin{tabular}{@{}lrr@{}}
\toprule
$N$ & $(N,N)$ float64 & with a 7-state shock \\
\midrule
200 & 0.3 MB & 2.2 MB \\
1{,}000 & 8 MB & 56 MB \\
10{,}000 & 800 MB & 5.6 GB \\
\bottomrule
\end{tabular}
\end{center}
\medskip
The arithmetic is still $O(N^2)$ per sweep. Vectorising changed the \emph{constant}, by a factor of about 60 --- it did not change the exponent.
\end{column}
\begin{column}{0.46\textwidth}
\begin{cautionbox}[title=\textbf{The honest limit}]\footnotesize
A $60\times$ constant is one and a half decades. The curse of dimensionality is many decades. \textbf{Vectorising buys you one more state variable, not five.}
\end{cautionbox}
\medskip
\begin{takeawaybox}\footnotesize
Which is why the next session attacks the exponent instead: exploiting monotonicity and concavity to avoid forming most of the matrix at all.
\end{takeawaybox}
\end{column}
\end{columns}
\end{frame}

%=============================================================================
\begin{frame}[fragile]{Measure Before You Optimise}
\begin{columns}[T]
\begin{column}{0.50\textwidth}
\begin{lstlisting}
%timeit m.bellman(V)
# 685 us per loop

import cProfile
cProfile.run("m.solve()", sort="cumtime")
\end{lstlisting}
\vspace{-1mm}
\small
Three questions, in this order:
\begin{enumerate}\footnotesize
  \item Is it \textbf{right}? Grade against the closed form first.
  \item Is it \textbf{too slow to use}? If not, stop.
  \item \textbf{Where} does the time go? Profile; never guess.
\end{enumerate}
\end{column}
\begin{column}{0.47\textwidth}
\begin{cautionbox}[title=\textbf{The two classic wastes}]\footnotesize
\textbf{Recomputing invariants} inside the loop --- the payoff matrix, the transition matrix, anything built from parameters alone.\\[1mm]
\textbf{Optimising the wrong 5\%.} In this solver almost all the time is in one line; the rest is rounding error in your afternoon.
\end{cautionbox}
\medskip
\begin{examplebox}\footnotesize
And record what you measured. A speed-up claim with no timing next to it is the same species of unsupported number as an accuracy claim with no benchmark.
\end{examplebox}
\end{column}
\end{columns}
\end{frame}

%=============================================================================
\begin{frame}{Takeaways}
\begin{takeawaybox}
\begin{itemize}\small
  \item Python is slow only when \textbf{Python} does the arithmetic. The same Bellman update, written as one array expression, ran about \textbf{65 times faster} and returned a bit-for-bit identical answer.
  \item \textbf{Broadcasting} is the whole trick: give every index of the mathematics its own axis, let the shapes stretch, then reduce along the axis you were maximising over.
  \item It almost never errors. A wrong axis gives you a plausible array of the wrong shape, so \textbf{assert the shape} and check against a known answer.
  \item A model is state plus behaviour, so make it a \textbf{class} --- and keep the closed-form benchmark inside it, where it cannot be evaluated at the wrong parameters.
  \item Vectorising changes the constant, not the exponent. Sixty times is one extra state variable, and that is all.
\end{itemize}
\end{takeawaybox}
\end{frame}

%=============================================================================
\begin{frame}{Bridge: The Derivative You Never Have to Code}
\begin{columns}[T]
\begin{column}{0.53\textwidth}
\small
NumPy gave us the arithmetic. It does not give us derivatives --- and from next week on, almost everything wants one: Newton's method needs a Jacobian, perturbation needs high-order derivatives at the steady state, and every neural network is trained by differentiating a loss.
\medskip

The old answer was finite differences, with a step size you have to tune against exactly the round-off of 3.A1. The modern answer is to let the library track the arithmetic as you do it, and hand you the exact derivative afterwards.
\end{column}
\begin{column}{0.44\textwidth}
\begin{conceptbox}[title=\textbf{Next (3.B3)}]\small
PyTorch from the ground up: tensors, automatic differentiation, \texttt{nn.Module}, and the training loop.\\[1mm]
Then how to work with an AI assistant on all of this --- and the one responsibility you cannot delegate.
\end{conceptbox}
\end{column}
\end{columns}
\end{frame}

\end{document}
